% GENERATED FILE. Edit canonical lessons, then run npm run generate:books.
% canonical-source-hash: fnv1a64:6b2e4ee3e5cca50d
% canonical-lessons: TA-C231-mayanku, TA-C231-vatu, TA-C231-eccari, TA-C231-kencu, TA-C231-camali

\chapter{To faint, To wilt, To warn, To beg, To manage}
\label{ch:ta-to-faint-to-wilt-to-warn-to-beg-to-manage}
\hlchaptermodality{231}

\begin{chapteropening}
\textbf{By the end of this chapter:} \emph{I can say to faint, to wilt, to warn, to beg and to manage.}

Complete the last lesson of chapter 231: I can say to faint, to wilt, to warn, to beg and to manage.
\end{chapteropening}

\section[\texorpdfstring{\ta{மயங்கு} (mayaṅku)}{மயங்கு (mayaṅku)}]{\ta{மயங்கு} (mayaṅku) — to faint, to pass out}

\label{lesson:TA-C231-mayanku}



\begin{quote}
Before the new one: say the Tamil for to download, then the Tamil for to upload.
\end{quote}

\subsection*{You'll want to know: \ta{மயங்கு}}
\textbf{\ta{மயங்கு}} — \emph{mayaṅku} — \textquotedblleft{}to faint, to pass out\textquotedblright{}.

\begin{grammarlens}[title={What you've built}]
One more word for this chapter.
\end{grammarlens}

\subsection*{Your turn}
\begin{itemize}
\raggedright
  \item \emph{Say it:} \emph{mayaṅku}
  \item \emph{Say it:} \emph{mayaṅku}, once more
  \item \emph{Say it:} say \emph{mayaṅku}
  \item \emph{Recall:} say the Tamil for to download, then the Tamil for to upload, then say \emph{mayaṅku} again
\end{itemize}

\begin{tcolorbox}[breakable,colback=teal!4,colframe=teal!35!black,title={Before you move on}]
What does \ta{மயங்கு} mean? (\textquotedblleft{}To faint, to pass out\textquotedblright{}.) Say it once more.
\end{tcolorbox}
\section[\texorpdfstring{\ta{வாடு} (vāṭu)}{வாடு (vāṭu)}]{\ta{வாடு} (vāṭu) — to wilt, to wither}

\label{lesson:TA-C231-vatu}



\begin{quote}
Before the new one: say the Tamil for to upload, then the Tamil for to faint.
\end{quote}

\subsection*{You'll want to know: \ta{வாடு}}
\textbf{\ta{வாடு}} — \emph{vāṭu} — \textquotedblleft{}to wilt, to wither\textquotedblright{}.

\begin{grammarlens}[title={What you've built}]
One more word for this chapter.
\end{grammarlens}

\subsection*{Your turn}
\begin{itemize}
\raggedright
  \item \emph{Say it:} \emph{vāṭu}
  \item \emph{Say it:} \emph{vāṭu}, once more
  \item \emph{Say it:} say \emph{vāṭu}
  \item \emph{Recall:} say the Tamil for to upload, then the Tamil for to faint, then say \emph{vāṭu} again
\end{itemize}

\begin{tcolorbox}[breakable,colback=teal!4,colframe=teal!35!black,title={Before you move on}]
What does \ta{வாடு} mean? (\textquotedblleft{}To wilt, to wither\textquotedblright{}.) Say it once more.
\end{tcolorbox}
\section[\texorpdfstring{\ta{எச்சரி} (eccari)}{எச்சரி (eccari)}]{\ta{எச்சரி} (eccari) — to warn}

\label{lesson:TA-C231-eccari}



\begin{quote}
Before the new one: say the Tamil for to faint, then the Tamil for to wilt.
\end{quote}

\subsection*{You'll want to know: \ta{எச்சரி}}
\textbf{\ta{எச்சரி}} — \emph{eccari} — \textquotedblleft{}to warn\textquotedblright{}.

\begin{grammarlens}[title={What you've built}]
One more word for this chapter.
\end{grammarlens}

\subsection*{Your turn}
\begin{itemize}
\raggedright
  \item \emph{Say it:} \emph{eccari}
  \item \emph{Say it:} \emph{eccari}, once more
  \item \emph{Say it:} say \emph{eccari}
  \item \emph{Recall:} say the Tamil for to faint, then the Tamil for to wilt, then say \emph{eccari} again
\end{itemize}

\begin{tcolorbox}[breakable,colback=teal!4,colframe=teal!35!black,title={Before you move on}]
What does \ta{எச்சரி} mean? (\textquotedblleft{}To warn\textquotedblright{}.) Say it once more.
\end{tcolorbox}
\section[\texorpdfstring{\ta{கெஞ்சு} (keñcu)}{கெஞ்சு (keñcu)}]{\ta{கெஞ்சு} (keñcu) — to beg, to plead}

\label{lesson:TA-C231-kencu}



\begin{quote}
Before the new one: say the Tamil for to wilt, then the Tamil for to warn.
\end{quote}

\subsection*{You'll want to know: \ta{கெஞ்சு}}
\textbf{\ta{கெஞ்சு}} — \emph{keñcu} — \textquotedblleft{}to beg, to plead\textquotedblright{}.

\begin{grammarlens}[title={What you've built}]
One more word for this chapter.
\end{grammarlens}

\subsection*{Your turn}
\begin{itemize}
\raggedright
  \item \emph{Say it:} \emph{keñcu}
  \item \emph{Say it:} \emph{keñcu}, once more
  \item \emph{Say it:} say \emph{keñcu}
  \item \emph{Recall:} say the Tamil for to wilt, then the Tamil for to warn, then say \emph{keñcu} again
\end{itemize}

\begin{tcolorbox}[breakable,colback=teal!4,colframe=teal!35!black,title={Before you move on}]
What does \ta{கெஞ்சு} mean? (\textquotedblleft{}To beg, to plead\textquotedblright{}.) Say it once more.
\end{tcolorbox}
\section[\texorpdfstring{\ta{சமாளி} (camāḷi)}{சமாளி (camāḷi)}]{\ta{சமாளி} (camāḷi) — to manage, to cope with}

\label{lesson:TA-C231-camali}



\begin{quote}
Before the new one: say the Tamil for to warn, then the Tamil for to beg.
\end{quote}

\subsection*{You'll want to know: \ta{சமாளி}}
\textbf{\ta{சமாளி}} — \emph{camāḷi} — \textquotedblleft{}to manage, to cope with\textquotedblright{}.

\begin{grammarlens}[title={What you've built}]
One more word for this chapter.
\end{grammarlens}

\subsection*{Your turn}
\begin{itemize}
\raggedright
  \item \emph{Say it:} \emph{camāḷi}
  \item \emph{Say it:} \emph{camāḷi}, once more
  \item \emph{Say it:} say \emph{camāḷi}
  \item \emph{Recall:} say the Tamil for to warn, then the Tamil for to beg, then say \emph{camāḷi} again
\end{itemize}

\begin{tcolorbox}[breakable,colback=teal!4,colframe=teal!35!black,title={Before you move on}]
What does \ta{சமாளி} mean? (\textquotedblleft{}To manage, to cope with\textquotedblright{}.) Say it once more.
\end{tcolorbox}
